\documentclass[11pt,a4paper]{article}
\usepackage[margin=24mm]{geometry}
\usepackage{amsmath,amssymb,booktabs,array}
\usepackage[hidelinks]{hyperref}
\setlength{\parindent}{0pt}
\setlength{\parskip}{3.8pt}
\setlength{\emergencystretch}{2em}
\title{How much traffic order is needed for useful\\urban signal-control gradients?}
\author{Urban methodology review: compact proposal}
\date{27 October 2026}
\begin{document}
\maketitle
\begin{abstract}
An empty turning bay cannot serve traffic that cannot reach its opening.
The same aggregate movement totals can therefore support different responses
to a service change, depending on local lane order and physical access.
We propose a continuous-mass first-order traffic model with exact ordered
route metadata near an isolated junction. Known transfer capability constrains
new lane allocation; existing commitments remain fixed; current lane heads
share receiving supply while consuming elapsed average service. The model
retains timestep boundary histories and composes one-step responses into
clearance-complete total system time. Aggregate models may use fine-grained
observations through their own declared macro conversion. The scientific
test is whether retained order improves externally evaluated service
decisions after strong aggregate explanations are allowed to compete.
This is a methodology proposal, without empirical fidelity or gradient
usefulness results.
\end{abstract}

\textbf{Feedback 10 status.}
The starting bay scenario uses shared left/through permission. Only actual
lane admission binds; unadmitted mass has no frozen lane target. Its
movement-specific processing policy is agreed, but the physical access
interface with the common feeder remains open. Q4 is withdrawn; average
receiving A is retained with sequential head dependencies. Frozen waiting
targets below are an unselected earlier proposal, not the current model.
See the separate \href{urban_feedback10_response.md}{feedback 10 answer}.
This update answers and records decisions only.



\section{From a task definition to a traffic explanation}

The established urban review already fixes the service-allocation task,
whole-population objective and V0--V4 comparisons. What it did not fully
specify was the traffic calculation joining lane assignment, bay access,
average service, downstream receiving and residual traffic. Naming an
ordered-access operator is insufficient: two plausible operators can
conserve mass yet leave different queues and control responses.

The VSL-like completion standard is a presentation and definition standard:
a reader should reconstruct the proposed model and sensitivity test without
inventing a missing traffic rule. It supplies no urban traffic assumption.
The present proposal fills that definition gap while keeping methodological
approval, implementation correctness and empirical usefulness separate.

Consider a left-turn block behind through traffic in one common feeder.
The bay can have room while the turner remains inaccessible. More left
service initially has no accessible traffic to serve; more through service
can remove the front obstruction and later release bay demand. Conversely,
if the turner is first, it can reach the bay before the through block.
Aggregate counts alone do not identify these two ordered states. A macro
model with a legitimate history-based estimator may nevertheless infer the
relevant access condition; such success must count against the necessity
of retaining the complete order.

\section{The proposed traffic model}

Use a symmetric isolated four-leg junction with fixed routes and four
protected service groups: NS through/right, NS left, EW through/right and
EW left. Continuous fractions $g_p\geq0$, $\sum_pg_p=1$ allocate effective
average service. No executable within-cycle waveform is assumed. Each
movement's lane service rate is its saturation rate times its fraction,
capped by physical lane capacity.

Long-road propagation uses the inherited cumulative first-order model.
The local geometry contains a physical common feeder and an internal fork
into a finite bay and adjacent continuation. Every physical road section
owns storage exactly once. Routes persist as table reference and route
position; internal section crossings do not advance the junction route
position. Quantity stays continuous throughout.

Per-lane ordered ledgers preserve current movement and route identity.
The movement grid records mass composition along cumulative rank; rank is
not physical distance, and grid composition does not recreate lost order.
An assignment waiting upstream is a metadata view of the upstream mass,
not an additional physical queue.

\subsection*{Allocation uses known capability}

The joined proposal places new assignment at the eligible fork front,
where current branch capability is known. The earlier prototype assigned
before entry to the feeder; that location was not selected in the principle
discussion. This location is therefore an explicit proposal, not an
editorial equivalence claim.

Process homogeneous cohorts in exact source order. Each sees the current
queue-load estimate plus assignments made to earlier cohorts, including
their waiting mass. Ideal balancing fills the shortest legal queues as
far as current sending, receiving and physical access permit. It does not
move existing traffic to manufacture equal queues or use future signal
plans as driver knowledge.

For starting lane loads $h_b$, current entry bounds $u_b$, common sending
$s$ and cohort mass $q$, the proposed exact law is
\begin{equation}
 V=\min(q,s,\sum_bu_b),\qquad
 x_b=\min\{u_b,[\lambda-h_b]_+\},\quad \sum_bx_b=V.
 \label{eq:compact-allocation}
\end{equation}
Thus $h=(0.2,0)$, $u=(0.3,0.5)$ and $s=0.6$ give entry
$x=(0.2,0.4)$ and balanced resulting loads $(0.4,0.4)$.
The remaining $0.4$ stays physically upstream. This agrees with feedback 9:
known capability shapes the allocation from the outset.

The distinction matters more visibly if $u=(0,1)$ and $s=1$. Fresh allocation
then sends the whole unit to the accessible legal lane. First freezing an
unconstrained target $(0.4,0.6)$ would instead waste usable capacity.
If that vector really was committed earlier, however, today's capacity
does not authorize its cancellation.

Feedback 10 rejects frozen waiting targets. Only actual entry fixes new
lane occupancy; unadmitted mass remains unassigned with priority within its
movement. The allocator may inspect other movements. How that processing
maps to physically accessible paths in the common feeder remains open;
this answer-only update derives no replacement coupled law.

Feedback 10 selects shared left/through permission for the starting bay
scenario. Permissions remain scenario-dependent and matched across V2/V3
comparisons. The exclusive-target ablation remains a different illustration
and cannot be used unchanged for the selected shared-permission case.

\subsection*{Order determines elapsed service}

A lane serves its current head at the applicable average rate. When that
head clears, the next movement becomes eligible. A blocked zero-rate head
stops that lane; other lanes can continue. Limited receiving is shared
among \emph{currently ready heads}: green-weighted once per inbound leg,
then equally within a leg subject to caps, reusing unused shares. Whole-step
compensation that equalizes earlier total discharge is unselected.

For left mass $q_L$ before sufficient through mass, with nonbinding receiving,
\begin{equation}
 y_T=s_Tg_T\left(\Delta t-\frac{q_L}{s_Lg_L}\right),
 \qquad 0<\frac{q_L}{s_Lg_L}<\Delta t.
 \label{eq:compact-service}
\end{equation}
With unit timestep/saturation, $q_L=0.2$ and fractions $(0.4,0.6)$,
through discharge is $0.3$. Moving a small fraction from through to left
has through-discharge derivative $0.25$: the front block clears sooner.
This is a conditional hand response, not a complete-time benefit claim.

The node solver uses the proposed step-average receiving rate
$R_j/\Delta t$, where $R_j$ is the inherited one-step receiving mass.
Its internal head calculation produces timestep totals. Propagation still
uses the existing affine interpolation between cumulative boundary counts;
new event-resolved histories remain deferred. These two conventions are
explicitly distinguished, without introducing a new between-step research
target.

\section{One-step response and the complete objective}

All allowances use pre-state histories. Each admitted amount leaves one
physical owner and enters another; newly entering mass receives no
same-section discharge in that step. With the inherited compatible
delay/timestep restriction, sending cannot exceed resident mass and
receiving cannot exceed vacant storage. Paired transfers preserve total
and route mass; waiting metadata contributes no extra stock.

The objective includes desired-release waiting, moving traffic and queues
until the complete finite population leaves, including a fixed continuation
after the controlled interval. For the inherited timestep account,
\begin{equation}
 J_{\Delta t}=\Delta t\sum_k
 \left[\sum_o(Q_{o,k}+D_{o,k})+\sum_en_{e,k}\right].
 \label{eq:compact-tst}
\end{equation}
At a fixed starting state the immediate direct cost partial with respect
to current service can be zero. Service changes still alter next queues,
access, later spillback and eventual sink exits. An internal transfer
alone does not remove population from the system.

For fixed structural order and active constraints, the specification gives
the queue-fill, sharing, head-clearance and history derivatives. Compose
them through the actual full state update:
\begin{equation}
 X_{k+1}=F_k(X_k,u_k),\qquad
 P_{k+1}=F_{k,X}P_k+F_{k,u}\dot u_k.
 \label{eq:compact-chain}
\end{equation}
During fixed continuation $\dot u_k=0$, but the state response $P_k$
continues through clearance. No per-vehicle trajectory-gradient construction
is required. A second accounting expression, complete released population
minus cumulative sink exit, checks the same total-time response.

Use a common horizon qualified to clear the compared action neighborhood;
its numerical value is later work. A changing stopping index, a positive
residual cutoff or a nonclearing trajectory cannot silently substitute for
complete TST. Derivatives are conditional on unchanged active signatures;
head, capacity and allocation switches can create kinks.

\section{Comparisons that can challenge the hypothesis}

Retain the accepted progression: V0 checks matched neutral behavior; V1
isolates shared-lane order; V2 tests finite storage without deliberate access
obstruction; V3 combines storage and ordered access; V4 tests held-out
variation. Equal aggregate totals alone do not establish full informational
equivalence, and nominally identical lanes may be collapsed only when their
order, access, readiness and service states are correspondingly symmetric.

Allow each macro rival to receive the original fine time series through
its own stated converter or estimator. Disclose retained representation,
memory, update schedule and calibration data. Legitimate estimation of
blockage counts as aggregate success; a hidden persistent ordered ledger
counts as retained meso information. State-conditioned replay, updated
prediction and forecasting from a common prefix support different claims.
State which observations are fixed or regenerated for each derivative test.

The preferred native B1 remains
\href{https://arxiv.org/html/2310.12249v2}{Wei et al.'s moving/turn-queue formulation},
with effective-green input. Its causal short-link treatment must be qualified
rather than replacing it silently by an invalid shortcut. Preferred B2
remains lane-exit allocation with an explicitly sourced storage/access
adaptation. \href{https://www.sciencedirect.com/science/article/pii/S0191261525000578}
{Gong et al.'s exit-allocation theory} does not itself define
finite-bay entry; \href{https://attap.umd.edu/wp-content/uploads/2022/05/TRC522.pdf}
{Liu--Chang} provides a concrete physical-storage source
whose native timing must be reconciled. As agreed in feedback 8, exact
implementation-dependent native selection remains pending.

The specification also defines a controlled aggregate ablation on the same
physical graph. It replaces exact order by a homogeneous mixture closure,
while keeping action, S/R, sources and objective matched. Its closure is a
project proposal, not a claimed reproduction of either native rival.
It makes a core comparison reconstructible without excluding stronger
native explanations. Pure matched channels provide neutral equivalence
checks; mixed totals alone do not.
If its mixture closure fails a storage-only V2 check, that is a limitation
of this ablation, not an order-necessity result. The stronger storage-aware
rival must be qualified on V2 before a V3 difference carries that claim.

Independent hand/reference cases establish a correctness floor. Later
SUMO evaluation uses a declared fixed-cycle conversion of effective fractions,
common routes and desired releases, fixed calibration and native lane
behavior. Complete reference time includes insertion delay plus travel
duration for the whole master population; unfinished traffic is retained
or reported as noncompletion. Average service and executable signals are
not claimed identical.

First test prescribed service perturbations externally; then assess
downstream optimization. Keep forward fidelity, derivative correctness and
gradient usefulness as separate claims. No experiment has yet demonstrated
that retained order improves decisions.

\section{Feedback 10: what is settled and what remains}

The starting shared-lane scenario, actual-admission commitment, progressive
counting of earlier admissions and constant average receiving A are
recorded. New unadmitted mass carries no frozen lane target, so Q4's old
waiting-quota premise is absent and that question is withdrawn.

The remaining model issue is physical application of movement-specific
skipping: inspecting another movement is distinct from being able to pass
a blocked head in one physical lane. The earlier common-prefix allocation
gate and frozen-waiting construction are not a completed current model.
The \href{urban_feedback10_response.md}{separate response} explains the
statuses and the correction that A still has sequential head dependencies.
No replacement derivation or next task is opened here.

\clearpage
\appendix
\section{Historical questions and feedback 10 status}

The question text below is preserved as the material answered in feedback 10.
Status notes give the current outcome; older candidate recommendations are
historical and do not override those outcomes.

\textbf{Historical pre-feedback-10 wording follows.} Written equations illustrate a proposed completion. Continuous mass, exact order, progressive queue loads, capability-aware fresh allocation and current-head sharing remain agreed.

Permissions determine which allocation questions are active: a unique legal target removes discretionary lane choice. Question 4 applies only if binding multi-lane waiting targets exist; it is not a general gate for fresh allocation.

\subsection*{1. Which lanes may a left turner use?}

\textbf{Feedback 10: Resolved: B for the starting bay scenario; other scenarios may declare other permissions.}

A bay and an adjacent lane do not automatically give a left turner two alternatives. This matters because the agreed responsive allocator can react to queues only where the permission matrix offers a real choice. An exclusive-lane model can still study storage and ordered access, but it does not exercise adaptive left-turn allocation.

\textbf{Candidate A: left turns use only the bay; through/right use only the continuation.} Suppose the bay already holds 0.2 units and the other lane is empty. With enough room and transfer capacity, a new left-turn unit must all go to the bay. The resulting loads are 1.2 and 0. This gives a simple storage/access case. A turner cannot avoid a busy bay by choosing the adjacent lane.

\textbf{Candidate B: the bay is left-only, while the adjacent lane also permits left turns alongside through/right.} In the same unconstrained example, the new left unit splits 0.4/0.6, giving loads 0.6/0.6. Choice now responds to queues created by earlier control. The shared lane can also contain left and through heads in different orders, so its stop-line obstruction must use the already selected elapsed-service rule.

\textbf{I provisionally favor B if responsive lane allocation is intended to be active in the main bay experiment.} This revises the earlier draft's exclusive-lane illustration as a preferred main realization: that illustration left the selected allocator with no legal choice. A remains a valid simpler control, not a defective traffic model. V2 and V3 must use the same chosen permissions, as must their comparison models; otherwise a change in legal movements could be mistaken for an order effect. No extra control series is required merely by listing A.

\subsection*{2. Where should a new lane choice become binding?}

\textbf{Feedback 10: Resolved in principle: actual admission binds and later packets count earlier admissions. Physical placement is part of Q3.}

Where drivers first commit determines which queue snapshot shapes their choice. Keep the physical common feeder and bay opening inside the local zone under both candidates. The question concerns the decision location, not moving the bay opening or reopening the rejected mouth-aligned geometry.

\textbf{Candidate A: commit before entering the common feeder.} A left-turn unit observes bay/other loads 0.2/0 and receives targets 0.4/0.6 when both are legal. It carries those targets while travelling or waiting in the feeder. This describes an earlier lane-choice intention and fits the earlier prototype. Known capability at that entry boundary can constrain present transfer, but current branch room cannot be treated as guaranteed room when the unit later reaches the fork.

\textbf{Candidate B: make the new choice at the eligible fork front.} If the observed loads have become 0/0.5 by then, the same new unit instead receives targets 0.75/0.25, assuming ample entry capability. Earlier commitment would retain 0.4/0.6. Thus the locations can predict different lane populations even though both preserve valid commitments and use only current queue information.

\textbf{I provisionally favor B for this first joined model}, because the relevant current branch limits and the new allocation refer to the same physical entry attempt. Feedback 9 does not logically force this location: early choice is also possible if current transfer capability and future branch admission are distinguished honestly. If the intended traffic story is that drivers have already chosen their future lane before entering the feeder, A is the appropriate direction and the written model must reflect that. Under unique legal targets, this location alone does not create a lane-choice effect.

\subsection*{3. Does the unadmitted remainder receive a binding lane target now?}

\textbf{Feedback 10: Policy resolved: no frozen target; priority within movement; inspect other movements. Physical access/interface remains open.}

Known capability determines how much can enter now, but it does not by itself say what a waiting record commits the rest to. With loads 0.2/0, branch entry allowances 0.3/0.5 and common sending 0.6, both candidates first admit 0.2/0.4. They agree with feedback 9 on this present entry. Their difference concerns the 0.4 that remains physically upstream.

\textbf{Candidate A: assign the waiting remainder now, then keep that intention fixed.} After selecting entry, balance the remaining 0.4 as waiting targets 0.2/0.2. Later packets see that committed load, and the remainder cannot change destination merely because later queues or capacities differ. This is the exact waiting convention proposed in the specification, not an unconstrained target computed before current entry.

\textbf{Candidate B: leave the unadmitted remainder unassigned until a later entry attempt.} Only the mass actually admitted acquires a binding lane choice. If, at the next opportunity, only the second legal lane has room and remains preferable, all remaining 0.4 may enter it. Under A, only its already committed 0.2 can do so. B delays creation of a target; it does not cancel a genuinely valid old target.

\textbf{I provisionally retain A to respect the existing default that logical assignment creates fixed waiting commitments.} I am not asking whether already valid targets may be changed. The point to clarify is whether the capacity-aware allocator creates valid targets for the unadmitted remainder at all. Under B, later packets on a common feeder wait for that unassigned prefix to be processed; when processing resumes, they see earlier realized assignments. Choosing B would explicitly amend the logical-waiting-assignment default for new mass. It would not authorize retargeting traffic whose commitment is already valid.

\subsection*{4. If old multi-lane targets are binding, how should scarce entry be shared?}

\textbf{Feedback 10: Withdrawn: the selected empty-loaded model creates no such waiting targets. The alternatives below are not a current decision.}

\textbf{This question is conditional.} It applies to an already committed homogeneous remainder whose targets span multiple lanes. If the selected model leaves all unadmitted mass unassigned, this case does not arise in its ordinary empty-loaded development path and need not block it. The earlier mandatory A/B question for fresh allocation remains withdrawn.

Suppose a genuinely old remainder has targets 0.2/0.2, current lane loads 0/0.2, enough branch room, but common sending permits only 0.2 to enter. Fixed targets determine eventual destination totals; without a further rule they do not say which portions enter first.

\textbf{Candidate A: advance the remaining targets proportionally, subject to caps and reuse.} Admit 0.1/0.1, leaving lane loads 0.1/0.3. This treats the target shares as progressing at the same fractional pace when no individual cap intervenes. It adds that interpretation of partial progress to the fixed-total commitment.

\textbf{Candidate B: balance current queue loads within the fixed target caps.} Admit 0.2/0, leaving loads 0.2/0.2. The remaining target is then 0/0.2. No destination total is changed: the second share waits longer, while more of the first share enters now. This treats old targets as fixed total quotas rather than a requirement of proportional progress.

\textbf{I provisionally favor B if this case is retained}, using the same queue-based criterion within the old caps. Neither candidate retargets the block or creates a lane-target sequence inside it. This is a choice about realizing a valid old commitment under a new limit, not a return to fixing fresh targets without known capability.

\subsection*{5. Does downstream receiving act as an average rate or a consumable timestep allowance?}

\textbf{Feedback 10: Resolved: retain A. Head clearance and current-head competition still have sequential dependencies.}

The inherited propagation model supplies an amount that the downstream section may receive during the update. The elapsed head solver also needs to know how that amount limits its internal service calculation. Two conservative interpretations can leave different lane queues at the next boundary. This is why the choice belongs to the present timestep model.

Consider a one-unit timestep and an exit that can receive 0.6 units, although its physical rate capacity is 1. Lane A has a ready head for that exit throughout. Lane B first serves an existing 0.5-unit head to another exit at rate 1; its next, already resident head then joins A halfway through the step. Both lanes have head-rate cap 1 and equal green weights when competing for the common exit. Sending and the other exit are nonbinding.

\textbf{Candidate A: spread the receiving allowance as a constant average rate.} The exit offers rate 0.6 for the whole step. A alone receives 0.3 during the first half; each lane receives another 0.15 during the second. Final common-exit discharges are 0.45/0.15. This is the receiving convention in the current equations. It puts receiving on the same step-average-rate basis as effective service, but stretches the allowance across the interval.

\textbf{Candidate B: serve at the physical receiving rate until the timestep allowance is spent.} A receives 0.5 in the first half. The remaining 0.1 is then shared by the two current heads, giving final discharges 0.55/0.05. The solver respects physical rate 1 and total receiving 0.6; it treats that allowed amount as available to current service rather than reserving a uniform fraction for later in the step.

\textbf{I provisionally retain A for the current average-rate model}, while making its effect on the advantage of an earlier head explicit. B is the better match if the intended meaning is a pool of currently usable downstream room. Neither is deduced from the endpoint allowance alone. Both retain current-ready-head sharing, effective signal rates and the existing boundary-history interpolation; neither equalizes final lane totals or introduces actual signal phases. The different final discharges make this a boundary-state choice, not a request for an intermediate-time gradient or the deferred event-history upgrade.

\subsection*{What can remain later}

The homogeneous-mixture ablation is optional. Its current formula covers exclusive targets; shared permissions require a matching comparison definition. It cannot replace stronger native rivals.

Numerical calibration, native-rival qualification, horizon values and reference-program settings remain later work. Bookkeeping and clearance conventions need new discussion only if they conflict with the chosen traffic behavior.

\textbf{Current topic: feedback 10 answers only. Physical access under movement-specific skipping remains open; no replacement coupled model is derived.}

\end{document}
